\documentclass{article}
\usepackage{amsmath,amssymb,amsthm}
\usepackage{graphicx,booktabs,xspace}
\usepackage[numbers]{natbib}
\newtheorem{thm}{Theorem}
\newtheorem{lem}[thm]{Lemma}
\newcommand{\sys}{\textsc{Breeze}\xspace}

\title{Energy-Aware Scheduling of Sensor Uploads}
\author{A. N. Author}
\date{}

\begin{document}
\maketitle

\section{Introduction}

In today's rapidly evolving landscape of urban sensing, efficient data
transmission plays a pivotal role in the longevity of battery-powered sensor
networks. Notably, the radio often accounts for the majority of a sensor's
energy budget---a fact that is frequently overlooked. This paper delves into
the intricate trade-off between upload frequency and data freshness.
Overall, this highlights the need for a comprehensive and principled
approach.

Most low-cost sensors upload each reading as soon as it is taken. This is
simple, but the radio dominates the energy budget; waking it up costs as
much as several hundred readings \citep{anastasi2009energy}. Batching the
uploads allows to save most of that energy, at the price of data that
arrives late. In what follows, we ask how long a sensor can wait before the
delay starts to matter.

\section{Method}

We model the value of a reading as decreasing with its age. A reading of age
$a$ has value $v(a) = e^{-a / \tau}$, where $\tau$ is set by the
application; a dashboard that people check every hour has a much larger
$\tau$ than an alarm. The sensor uploads a batch when the value it would
lose by waiting longer exceeds the energy it would save.

\begin{thm}\label{thm:batch}
  Under the model above, the optimal batch size grows with the square root of $\tau$.
\end{thm}

Put differently, an application that tolerates four times the delay should
wait for twice as many readings. We stress that this rule needs no knowledge
of the radio beyond the energy cost of one wake-up, which the manufacturer
usually reports.

\section{Experiments}

We ran the scheduler on 12 sensors for six weeks, next to 12 sensors that
upload every reading. The scheduler sends less samples per upload than a
fixed batch of 50 would, but it uploads them when they are still useful.
Battery life went from 21 to 47 days, and the median delay of a reading was
9 minutes.

Our findings underscore the importance of adaptive batching. Moreover, the
proposed scheduler seamlessly integrates with existing firmware, paving the
way for more sustainable deployments. In essence, it is not just an energy
optimization but a holistic rethinking of how sensors communicate.

\section{Conclusion}

We presented a scheduler that batches uploads by the value of the data. It
more than doubles battery life in our deployment, and it needs one parameter
that the application already knows.

\end{document}
